\documentclass[10pt,a4paper]{amsart}
\usepackage{amsmath,amssymb,amsthm,graphicx}
\usepackage{verbatim}
\newenvironment{reference}{\comment}{\endcomment}
%\newenvironment{reference}{}{}
\newenvironment{slogan}{\comment}{\endcomment}
\newenvironment{history}{\comment}{\endcomment}

% For commutative diagrams we use Xy-pic
\usepackage[all]{xy}

% We use 2cell for 2-commutative diagrams.
\xyoption{2cell}
% Theorem environments.
%
\theoremstyle{plain}
\newtheorem{theorem}[subsection]{Theorem}
\newtheorem{proposition}[subsection]{Proposition}
\newtheorem{lemma}[subsection]{Lemma}

\theoremstyle{definition}
\newtheorem{definition}[subsection]{Definition}
\newtheorem{example}[subsection]{Example}
\newtheorem{exercise}[subsection]{Exercise}
\newtheorem{situation}[subsection]{Situation}

\theoremstyle{remark}
\newtheorem{remark}[subsection]{Remark}
\newtheorem{remarks}[subsection]{Remarks}

\numberwithin{equation}{subsection}

% Macros
%
\def\lim{\mathop{\mathrm{lim}}\nolimits}
\def\colim{\mathop{\mathrm{colim}}\nolimits}
\def\Spec{\mathop{\mathrm{Spec}}}
\def\Hom{\mathop{\mathrm{Hom}}\nolimits}
\def\Ext{\mathop{\mathrm{Ext}}\nolimits}
\def\SheafHom{\mathop{\mathcal{H}\!\mathit{om}}\nolimits}
\def\SheafExt{\mathop{\mathcal{E}\!\mathit{xt}}\nolimits}
\def\Sch{\mathit{Sch}}
\def\Mor{\mathop{\mathrm{Mor}}\nolimits}
\def\Ob{\mathop{\mathrm{Ob}}\nolimits}
\def\Sh{\mathop{\mathit{Sh}}\nolimits}
\def\NL{\mathop{N\!L}\nolimits}
\def\CH{\mathop{\mathrm{CH}}\nolimits}
\def\proetale{{pro\text{-}\acute{e}tale}}
\def\etale{{\acute{e}tale}}
\def\QCoh{\mathit{QCoh}}
\def\Ker{\mathop{\mathrm{Ker}}}
\def\Im{\mathop{\mathrm{Im}}}
\def\Coker{\mathop{\mathrm{Coker}}}
\def\Coim{\mathop{\mathrm{Coim}}}

% Boxtimes
%
\DeclareMathSymbol{\boxtimes}{\mathbin}{AMSa}{"02}

%
% Macros for moduli stacks/spaces
%
\def\QCohstack{\mathcal{QC}\!\mathit{oh}}
\def\Cohstack{\mathcal{C}\!\mathit{oh}}
\def\Spacesstack{\mathcal{S}\!\mathit{paces}}
\def\Quotfunctor{\mathrm{Quot}}
\def\Hilbfunctor{\mathrm{Hilb}}
\def\Curvesstack{\mathcal{C}\!\mathit{urves}}
\def\Polarizedstack{\mathcal{P}\!\mathit{olarized}}
\def\Complexesstack{\mathcal{C}\!\mathit{omplexes}}
% \Pic is the operator that assigns to X its picard group, usage \Pic(X)
% \Picardstack_{X/B} denotes the Picard stack of X over B
% \Picardfunctor_{X/B} denotes the Picard functor of X over B
\def\Pic{\mathop{\mathrm{Pic}}\nolimits}
\def\Picardstack{\mathcal{P}\!\mathit{ic}}
\def\Picardfunctor{\mathrm{Pic}}
\def\Deformationcategory{\mathcal{D}\!\mathit{ef}}
\usepackage{cleveref}
\crefname{originallemma}{original lemma}{original lemmas}
\Crefname{originallemma}{Original lemma}{Original lemmas}
\makeatletter
\newlabel{chow-lemma-compose-pushforward}{{02R5}{source}{}{sourcepointer-1}{}}
\newlabel{chow-lemma-compose-pushforward@cref}{{[originallemma][0][]02R5}{[1][0][]source}{}{}{}}
\makeatother
\crefname{originallemma}{original lemma}{original lemmas}
\Crefname{originallemma}{Original lemma}{Original lemmas}
\makeatletter
\newlabel{chow-lemma-intersect-regularly-embedded}{{0FFE}{source}{}{sourcepointer-2}{}}
\newlabel{chow-lemma-intersect-regularly-embedded@cref}{{[originallemma][0][]0FFE}{[1][0][]source}{}{}{}}
\makeatother
\crefname{originallemma}{original lemma}{original lemmas}
\Crefname{originallemma}{Original lemma}{Original lemmas}
\makeatletter
\newlabel{chow-lemma-lci-gysin-composition}{{0FF5}{source}{}{sourcepointer-3}{}}
\newlabel{chow-lemma-lci-gysin-composition@cref}{{[originallemma][0][]0FF5}{[1][0][]source}{}{}{}}
\makeatother
\crefname{originallemma}{original lemma}{original lemmas}
\Crefname{originallemma}{Original lemma}{Original lemmas}
\makeatletter
\newlabel{chow-lemma-lci-gysin-flat}{{0FF4}{source}{}{sourcepointer-4}{}}
\newlabel{chow-lemma-lci-gysin-flat@cref}{{[originallemma][0][]0FF4}{[1][0][]source}{}{}{}}
\makeatother
\crefname{originalreference}{original reference}{original references}
\Crefname{originalreference}{Original reference}{Original references}
\makeatletter
\newlabel{example-graph-correspondence}{{0FG1}{source}{}{sourcepointer-5}{}}
\newlabel{example-graph-correspondence@cref}{{[originalreference][0][]0FG1}{[1][0][]source}{}{}{}}
\makeatother
\crefname{originalreference}{original reference}{original references}
\Crefname{originalreference}{Original reference}{Original references}
\makeatletter
\newlabel{remark-transpose}{{0FG4}{source}{}{sourcepointer-6}{}}
\newlabel{remark-transpose@cref}{{[originalreference][0][]0FG4}{[1][0][]source}{}{}{}}
\makeatother
\title{The correspondence of a graph acting on cycles}
\author{The Stacks Project Authors}
\date{stacks-project:weil.tex:local-snapshot-20261012; GNU FDL 1.2 or later}
\begin{document}
\maketitle
\noindent\textit{Editorial note.} All delivered mathematical text is the author's own, copied from the pinned Stacks Project chapter file weil.tex; no statement or proof step was rewritten, reordered or removed.\par\medskip
\noindent\textit{Editorial note.} The counted proof is the author's complete argument: Prefilter hit re-read per the 2026-10-14 12:4x THIN-criterion patch (pb>=800 / pw>=100): the proof body was read in full at source and its main line is complete; the flagged phrase is incidental and is quoted verbatim below. Verbatim, proof line 713/714: 'Proof of (3). Proved in exactly the same manner as above.'  -  items (1) and (2) of the statement are derived explicitly in the span and item (3) is their transpose variant obtained by interchanging the two projections, so the main line of the argument is present.\par\medskip
\noindent\textit{Editorial note.} The label lemma-functor-and-cycles names the counted statement, and the proof environment immediately following it in the source is delivered verbatim as its proof.\par\medskip
\noindent\textit{Editorial note.} The counted statement is carried by the author's own lemma environment in the chapter source.\par\medskip
\noindent\textit{Editorial note.} Prefilter hit re-read per the 2026-10-14 12:4x THIN-criterion patch: proof body (proof\_lines 683-715, pb=981, pw=122) read in full; items (1) and (2) derived explicitly, item (3) a transpose repetition.\par\medskip
\noindent\textit{Editorial note.} Editorial note (verbatim): 'Proof of (3). Proved in exactly the same manner as above.'\par\medskip
\noindent\textit{Editorial note.} Pre-open gate: python operations/agent-authoring/tools/precard\_gate.py --source weil.tex --label lemma-functor-and-cycles --start 660 --end 715 --local-id 28166 returned RESULT=OK. Pre-assembly tag self-check: canonical tag resolved and not present in reference/public-items.jsonl.\par\medskip
\noindent\textit{Editorial note.} The statement and proof are copied from the pinned chapter file; the cited inputs are kept as external pointers with their canonical Stacks tags.\par\medskip
\noindent\textit{Allowed standard prerequisites (editorial).}\begin{itemize}
\item Intersection with a regularly embedded subscheme and composition of Gysin maps. Chow Homology, the lemmas identifying the intersection with a regularly embedded subscheme through the Gysin map, computing the Gysin map for flat morphisms, and composing Gysin maps. Rewrites the intersection of the graph with a pulled back cycle in terms of the Gysin map along the graph.
\item Composition of proper pushforwards. Chow Homology, the lemma that proper pushforward is compatible with composition. Identifies the pushforward along the projection composed with the graph with the identity pushforward.
\end{itemize}\medskip
\ifdefined\hypertarget\hypertarget{sourcepointer-1}{}\fi
\noindent\textit{Source reference.} Original lemma 02R5: sources/stacks-chow/tags.txt, source lines 3612--3612. Prerequisite pointer invoked at proof line(s) 702; it names a standard fact of a sibling chapter that the counted argument uses, and every step of the counted proof built on it is retained verbatim in the delivered body.\par\medskip
\ifdefined\hypertarget\hypertarget{sourcepointer-2}{}\fi
\noindent\textit{Source reference.} Original lemma 0FFE: sources/stacks-chow/tags.txt, source lines 19284--19284. Prerequisite pointer invoked at proof line(s) 695; it names a standard fact of a sibling chapter that the counted argument uses, and every step of the counted proof built on it is retained verbatim in the delivered body.\par\medskip
\ifdefined\hypertarget\hypertarget{sourcepointer-3}{}\fi
\noindent\textit{Source reference.} Original lemma 0FF5: sources/stacks-chow/tags.txt, source lines 19275--19275. Prerequisite pointer invoked at proof line(s) 699; it names a standard fact of a sibling chapter that the counted argument uses, and every step of the counted proof built on it is retained verbatim in the delivered body.\par\medskip
\ifdefined\hypertarget\hypertarget{sourcepointer-4}{}\fi
\noindent\textit{Source reference.} Original lemma 0FF4: sources/stacks-chow/tags.txt, source lines 19274--19274. Prerequisite pointer invoked at proof line(s) 697; it names a standard fact of a sibling chapter that the counted argument uses, and every step of the counted proof built on it is retained verbatim in the delivered body.\par\medskip
\ifdefined\hypertarget\hypertarget{sourcepointer-5}{}\fi
\noindent\textit{Source reference.} Original reference 0FG1: sources/stacks-weil/pinned-original-source.tex, source lines 534--534. Prerequisite pointer invoked at proof line(s) 664; it names a standard fact of the same chapter that the counted argument uses, and every step of the counted proof built on it is retained verbatim in the delivered body.\par\medskip
\ifdefined\hypertarget\hypertarget{sourcepointer-6}{}\fi
\noindent\textit{Source reference.} Original reference 0FG4: sources/stacks-weil/pinned-original-source.tex, source lines 645--645. Prerequisite pointer invoked at proof line(s) 674; it names a standard fact of the same chapter that the counted argument uses, and every step of the counted proof built on it is retained verbatim in the delivered body.\par\medskip
\begin{lemma}
\label{lemma-functor-and-cycles}
Let $f : Y \to X$ be a morphism of smooth projective schemes over $k$.
Let $[\Gamma_f] \in \text{Corr}^0(X, Y)$ be as in
Example \ref{example-graph-correspondence}. Then
\begin{enumerate}
\item pushforward of cycles by the correspondence $[\Gamma_f]$
agrees with the Gysin map $f^! : \CH^*(X) \to \CH^*(Y)$,
\item pullback of cycles by the correspondence $[\Gamma_f]$
agrees with the pushforward map $f_* : \CH_*(Y) \to \CH_*(X)$,
\item if $X$ and $Y$ are equidimensional of dimensions $d$ and $e$,
then
\begin{enumerate}
\item pushforward of cycles by the correspondence
$[\Gamma_f^t]$ of Remark \ref{remark-transpose}
corresponds to pushforward of cycles by $f$, and
\item pullback of cycles by the correspondence
$[\Gamma_f^t]$ of Remark \ref{remark-transpose}
corresponds to the Gysin map $f^!$.
\end{enumerate}
\end{enumerate}
\end{lemma}

\begin{proof}
Proof of (1). Recall that
$[\Gamma_f]_*(\alpha) =
\text{pr}_{2, *}([\Gamma_f] \cdot \text{pr}_1^*\alpha)$.
We have
$$
[\Gamma_f] \cdot \text{pr}_1^*\alpha =
(f, 1)_*((f, 1)^! \text{pr}_1^*\alpha) =
(f, 1)_*((f, 1)^! \text{pr}_1^!\alpha) =
(f, 1)_*(f^!\alpha)
$$
The first equality by Chow Homology, Lemma
\ref{chow-lemma-intersect-regularly-embedded}.
The second by 
Chow Homology, Lemma \ref{chow-lemma-lci-gysin-flat}.
The third because $\text{pr}_1 \circ (f, 1) = f$ and
Chow Homology, Lemma \ref{chow-lemma-lci-gysin-composition}.
Then we coclude because
$\text{pr}_{2, *} \circ (f, 1)_* = 1_*$ by
Chow Homology, Lemma \ref{chow-lemma-compose-pushforward}.

\medskip\noindent
Proof of (2). Recall that $[\Gamma_f]_*(\beta) =
\text{pr}_{1, *}([\Gamma_f] \cdot \text{pr}_2^*\beta)$.
Arguing exactly as above we have
$$
[\Gamma_f] \cdot \text{pr}_2^*\beta = (f, 1)_*\beta
$$
Thus the result follows as before.

\medskip\noindent
Proof of (3). Proved in exactly the same manner as above.
\end{proof}

\end{document}
